% Draft for ICAIL 2027 (paper or demo track). Compile on Overleaf with the ACM "acmart" template:
% New Project -> Upload Project -> this file + references.bib.
% Before submitting: check the call for the page limit and anonymity rules, and verify every reference.
\documentclass[sigconf,nonacm]{acmart}
\usepackage{booktabs}

% ---- Numbers from eval/results.json (held-out evaluation) ----
\newcommand{\nJudgments}{38{,}413}
\newcommand{\nMinedSC}{13{,}044}
\newcommand{\nMinedHC}{80{,}931}
\newcommand{\nAttest}{47{,}770}
\newcommand{\genHigh}{1.2\%}
\newcommand{\genReview}{6.5\%}
\newcommand{\genOk}{92.3\%}
\newcommand{\fakeCaseHigh}{47.9\%}\newcommand{\fakeCaseRev}{50.8\%}\newcommand{\fakeCaseMiss}{1.3\%}
\newcommand{\wrongCiteHigh}{84.5\%}\newcommand{\wrongCiteRev}{13.9\%}\newcommand{\wrongCiteMiss}{1.6\%}
\newcommand{\wrongYearHigh}{34.6\%}\newcommand{\wrongYearRev}{63.1\%}\newcommand{\wrongYearMiss}{2.3\%}
\newcommand{\fakeQuoteHigh}{99.0\%}\newcommand{\fakeQuoteRev}{0.9\%}\newcommand{\fakeQuoteMiss}{0.1\%}
\newcommand{\realQuoteHigh}{0.0\%}\newcommand{\realQuoteRev}{0.9\%}\newcommand{\realQuoteOk}{99.1\%}

\begin{document}

\title{Pramaan: Catching Fabricated Case Citations in Indian Legal Documents with Open Data}

\author{Vividh Yadav}
\affiliation{%
  \institution{Birla Institute of Technology, Mesra}
  \city{Ranchi}
  \country{India}}

\begin{abstract}
In July 2026 the Supreme Court of India set aside a tribunal's order because it rested on precedents
generated by artificial intelligence. Four of the six cases cited did not exist, and the two real ones
carried invented paragraphs. The Court declared zero tolerance for unverified AI-generated
precedents. Yet in India, verifying a citation usually requires a subscription to a proprietary law
report. We present \emph{Pramaan}, an open-source system that checks every case citation in a document
against open data alone. Pramaan combines four components:
(i) a grammar for Indian reporter and neutral citations;
(ii) the official record of \nJudgments{} Supreme Court judgments;
(iii) an \emph{open citation concordance}: \nAttest{} (citation, case) attestations mined from how
Supreme Court and High Court judgments cite earlier cases, which stands in for proprietary reporters;
(iv) n-gram verification of quoted passages.
We describe a failure mode we call \emph{attestation poisoning}: a judgment that condemns fake
citations also quotes them. We guard against it. On a held-out evaluation built from 2023--2026
judgments, Pramaan flags \genHigh{} of genuine citations as high-risk. It lets through
\fakeCaseMiss{} of invented cases, \wrongCiteMiss{} of misattributed citations, and \fakeQuoteMiss{}
of invented quotations. On the six authorities from the case the Court set aside, it flags all four
invented cases and the misattributed citation, and catches the invented paragraph. Pramaan also runs
entirely in a web browser, so confidential documents never leave the user's machine.
\end{abstract}

\keywords{legal citations, hallucination detection, AI and law, India, open data, verification}

\maketitle

\section{Introduction}

Large language models fabricate legal authority~\cite{dahl2024fictions,magesh2024hallucination}.
Indian courts have now faced the consequences.

In \emph{Pooja Ramesh Singh v.\ Jammu and Kashmir Bank Ltd.} (2026 INSC 668), the National Company
Law Tribunal (NCLT) had admitted an insolvency petition relying on six Supreme Court ``precedents''.
Four did not exist. The other two had real citations but invented paragraphs. The appellate tribunal
did not notice.

The Supreme Court set both orders aside. It held three things:
\begin{itemize}
\item citing AI-generated precedents without verification is professional misconduct;
\item a decision touched by ``an iota'' of such material is no decision in the eyes of the law;
\item the Bar Council of India must frame norms on the issue.
\end{itemize}
It also asked what this paper answers: ``Imagine the hardship of a situation in which the Court must
verify the authenticity of each judgment cited by an advocate.''

Verification is hard in India for a specific reason. The citations lawyers use, such as
\texttt{(2018) 16 SCC 228}, point into proprietary reporters. The open record of Supreme Court
judgments~\cite{scjudgments} contains the text and the official Supreme Court Reports (SCR)
citation, but no table that says which case \texttt{(2018) 16 SCC 228} is.

\paragraph{Contributions.}
\begin{enumerate}
\item An open grammar for Indian citations. It covers SCC (including Supplementary and series
volumes), SCR, AIR, SCC OnLine, INSC and High Court neutral citations, and seven further reporters. It
also groups parallel citations and recovers case names, pinpoints and quotations from free text.
\item An \emph{open citation concordance} mined from judgments themselves. Each time a judgment writes
``\emph{X v.\ Y}, (2020) 2 SCC 677 : [2019] 16 SCR 252'', it attests to which case a proprietary
citation belongs. The parallel SCR citation ties that case to the official record.
\item The identification of \emph{attestation poisoning} and a guard against it.
\item A held-out evaluation, built from 2023--2026 judgments, of the fabrication types the Supreme
Court itself described.
\item An open-source implementation: a command line, an API, and a fully client-side browser build.
\end{enumerate}

\section{Related Work}

\emph{eyecite}~\cite{cushman2021eyecite} extracts and resolves United States citations and inspired
our grammar. However, Indian reporters, neutral citations and case-title conventions differ
substantially.

Work on Indian legal NLP has focused on:
\begin{itemize}
\item named entities, including precedents~\cite{kalamkar2022ner};
\item benchmarks for understanding judgments~\cite{joshi2024iltur}.
\end{itemize}
Neither line of work verifies that a citation is real.

Studies of legal hallucination quantify how often models invent authority~\cite{dahl2024fictions}.
They also show that commercial research tools are not immune~\cite{magesh2024hallucination}.

Commercial citation checkers for Indian law exist. They are closed and tied to proprietary databases,
so their error rates cannot be audited.

\section{System}

\paragraph{Grammar.}
Fifteen patterns recognise reporter and neutral citations. Adjacent citations separated only by
punctuation are grouped as parallel citations of one case.

The case name is recovered by scanning backwards from the citation:
\begin{itemize}
\item to the last ``v.''/``vs.''/``versus'' marker;
\item then over name-like tokens, stopping at discourse words (``In'', ``See'', ``Law Cited''),
dates and paragraph breaks.
\end{itemize}
Quoted passages of at least eight words are attached to the authority that introduces them.

\paragraph{Official record.}
For each of the \nJudgments{} judgments (1950--2026), we index:
\begin{itemize}
\item parties and decision date;
\item SCR year, volume and page range (so a pinpoint citation resolves to its judgment);
\item neutral citation.
\end{itemize}

\paragraph{Open citation concordance.}
We extract every Supreme Court citation from \nMinedSC{} Supreme Court judgments (2010--2026) and
\nMinedHC{} High Court judgments.
\begin{itemize}
\item \textbf{Clustering.} Mentions of the same citation are clustered by party-name similarity.
\item \textbf{Votes.} A cluster's vote count is the number of distinct judgments that cite the case
that way.
\item \textbf{Resolution.} Each cluster is resolved to an official judgment directly (SCR or INSC),
through a parallel SCR citation, or by name and year.
\end{itemize}
The result covers cases absent from the official reports, such as those reported only in SCC.

\paragraph{Attestation poisoning.}
Our first concordance ``verified'' the NCLT's fabricated citations. The source was 2026 INSC 668
itself: it lists the fakes in order to condemn them.

More generally, any judgment discussing fabricated authority repeats it, often far from the words
that condemn it. We exclude:
\begin{itemize}
\item every judgment whose text repeatedly discusses fabricated, non-existent or AI-generated
authority;
\item any citation within 1,500 characters of such language.
\end{itemize}

\paragraph{Name matching.}
Indian case titles vary widely. Examples:
\begin{itemize}
\item abbreviations: UOI, CIT, SEB;
\item romanisation: Kesavananda/Keshavananda;
\item OCR damage: ``Kamataka'', ``Kera/a'';
\item truncation: ``Sarguja Transport Service v.\ STAT''.
\end{itemize}
We normalise abbreviations, reduce tokens to a transliteration skeleton, and score party pairs by
IDF-weighted containment. Both parties must agree: the pair score is pulled toward the weaker party.

\paragraph{Quote verification.}
A quotation is split at ellipses. We report the share of its word 5-grams found in the judgment, which
is robust to hyphenation, page furniture and layout, but not to invented sentences.

\paragraph{Verdicts.}
Each authority receives one of three verdicts: \emph{verified}, \emph{review}, or \emph{high risk}.
\begin{itemize}
\item \textbf{High risk} only when open evidence positively contradicts the citation:
\begin{itemize}
\item the citation belongs to another case;
\item a distinctive case name appears only in another year;
\item the quotation is absent;
\item a Supreme Court citation from a well-covered reporter and year matches nothing.
\end{itemize}
\item \textbf{Review} where open sources are thin: recent years, SCC OnLine (which also reports
unreported orders), or common names such as ``Ram Kumar v.\ State of U.P.''.
\end{itemize}

\section{Evaluation}

\paragraph{Held-out design.}
The concordance for evaluation is mined only from judgments decided up to 2022. Test citations are
drawn from judgments of 2023--2026, which the concordance never saw.

We generate fabrications following the Supreme Court's taxonomy in 2026 INSC 668:
\begin{itemize}
\item \emph{invented case}: the petitioner of one real case joined to the respondent of another;
\item \emph{misattributed citation}: a real case name with another real case's citation;
\item \emph{wrong year}: the year shifted by 3--10;
\item \emph{invented quotation}: a passage from a different judgment.
\end{itemize}
Table~\ref{tab:results} reports 1,000 cases per category.

\begin{table}[t]
\caption{Held-out results (1,000 cases per row). ``Missed'' includes references the parser did not
extract.}
\label{tab:results}
\begin{tabular}{lrrr}
\toprule
Test set & High & Review & Missed / Verified \\
\midrule
Genuine citations & \genHigh & \genReview & \genOk{} verified \\
Invented case & \fakeCaseHigh & \fakeCaseRev & \fakeCaseMiss{} missed \\
Misattributed citation & \wrongCiteHigh & \wrongCiteRev & \wrongCiteMiss{} missed \\
Wrong year & \wrongYearHigh & \wrongYearRev & \wrongYearMiss{} missed \\
Invented quotation & \fakeQuoteHigh & \fakeQuoteRev & \fakeQuoteMiss{} missed \\
Genuine quotation & \realQuoteHigh & \realQuoteRev & \realQuoteOk{} verified \\
\bottomrule
\end{tabular}
\end{table}

\paragraph{Error analysis.}
We audited all 12 genuine citations flagged high-risk by hand:
\begin{itemize}
\item 8 are real judgments absent from open sources (SCC reports cases the official SCR omits);
\item 3 come from text-extraction noise in the citing judgment;
\item 1 carries a year that disagrees with the official record, a typo or OCR error at the source.
\end{itemize}
Adding 80{,}931 High Court judgments to the concordance reduced false alarms from 1.6\% to 1.2\%,
which suggests that coverage, not the method, is the main limit.

\paragraph{The NCLT's six authorities.}
On the six citations of 2026 INSC 668, Pramaan:
\begin{itemize}
\item flags all four non-existent cases;
\item marks the misattributed SCC OnLine citation for review;
\item verifies the genuine \emph{Canara Bank} citation, and catches an invented paragraph attributed
to it.
\end{itemize}
It errs on \emph{Everest Kento Cylinders}. The citation is real but absent from open sources, so it
is reported as not found. The citation was nonetheless problematic: its quoted paragraph was invented.

\section{Deployment}

Pramaan is a Python package with no required dependencies. It offers:
\begin{itemize}
\item a command line that exits non-zero when a document contains a high-risk citation, for use in
filing pipelines;
\item a JSON API with per-key quotas;
\item a static site that runs the same code in the browser via WebAssembly (Pyodide) with a 5\,MB
compressed index.
\end{itemize}
In the browser build, the document is never uploaded. Only public judgment PDFs are fetched, to check
quotations. This matters for advocates bound by client confidentiality.

\section{Limitations and Ethics}

``Not found'' means not found in \emph{open} sources. A false accusation of fabrication is a serious
harm, so Pramaan:
\begin{itemize}
\item reserves high risk for positive contradictions;
\item states its coverage in every finding;
\item presents findings as evidence to check, not rulings.
\end{itemize}
Coverage is limited to Supreme Court authorities. Extending it to High Courts, whose open dataset
holds 17.8 million judgments, is future work. Quotations can only be checked where judgment text is
open.

\section{Conclusion}

Open judgments, read carefully, can verify proprietary citations. A tool built on them can be audited,
run privately, and adopted by courts and the Bar at no cost. Code, data and evaluation are available at
\url{https://github.com/VividhDesign/pramaan}.

\bibliographystyle{ACM-Reference-Format}
\bibliography{references}

\end{document}
